% Supplementary File 1. Created 2026-08-09 by Phase 1 of papers/1_method/08_length_plan.md.
% NOT input by elife_paper.tex: eLife supplementary files are uploaded as separate documents and
% are labelled "Supplementary File 1, 2, ..." in the text. Everything here was moved verbatim out
% of Methods; nothing in the manuscript's argument depends on it, and Methods carries a one-line
% pointer at each of the four places it left.
%
% eLife's rule that bounds what may live here: "Do not bury detailed methods in supplements."
% What moved is not method, it is inventory: which cells exist on disk, which files each stage
% writes, the cluster's resource settings, and the optimiser's numeric schedule. The method itself,
% including the convergence criterion that actually stops the fits, stayed in Methods.

\documentclass[11pt]{article}
\usepackage[margin=1in]{geometry}
\usepackage{amsmath,amssymb}
\usepackage{bm}
\title{Supplementary File 1: the measured grid, the reduction stages, and the production environment}
\author{}
\date{}
\begin{document}
\maketitle

\section*{S1.1\quad The measured grid, cell by cell}

Across the three directories, \texttt{IR} covers
$N_{\mathrm{ch}} \in \{5, 10, 20, 50, 100, 200, 500, 1000, 2000, 5000, 10000\}$ and
all twelve levels; NR covers
$N_{\mathrm{ch}} \in \{10, 100, 1000, 10000\}$ and R covers those and additionally
$5$, $20$ and $50$ at $\widetilde{S} = 0.01$, with the highest levels run only at the larger
channel counts and reaching $\widetilde{S} = 10^6$ at $N_{\mathrm{ch}} = 10^4$; MR and
VR cover the same channel numbers at $\widetilde{S} \in \{0.01, 0.1, 1\}$; and the least-squares arm
covers $30$ cells on $N_{\mathrm{ch}} \in \{10, 100, 1000, 10000\}$, with $\widetilde{S}$ from
$0.01$ up to $10^3$, $10^4$, $10^5$ and $10^6$ respectively.
% src: projects/eLife_2025/figures/data/1c2ae6f + 87889e6 + 0ffbda7 (per-algorithm nch and noise sets, nsim_10000 files)

\section*{S1.2\quad The five reduction stages and what each writes}

Each grid cell is reduced through five stages: the per-group MLE cloud; the pooled
joint fit $\theta_{\mathrm{pool}}$; the score at both anchors, whose
across-replicate covariance is $J$ and whose per-interval Gaussian terms sum to the
anchor $\bar G$; the diagnostic battery at each anchor; and the
empirical-versus-theoretical distortion capstone, anchored at
$\theta_{\mathrm{pool}}$. Each cell writes five comma-separated-value families (the
estimate cloud, the pooled fit, the two anchor batteries, and the capstone). The
scripts number these stages one, two, four, five and six; the gap at three is the
numerical-Fisher stage, which the Gaussian anchor removes and which the numbering
keeps visible. On the grid-production lane the cloud stage and the capstone are
dropped, because no map reads them there; the cells that need a cloud are produced
by the separate script that keeps that stage.

\section*{S1.3\quad The production environment}

The production batteries ran under SLURM at $32$ CPUs per task and $48\;\mathrm{GB}$
of memory ($96\;\mathrm{GB}$ for the parameter-recovery lane), with the linear-algebra
back end pinned to a single thread so that OpenMP parallelises the per-simulation and
bootstrap loops. The environment variable
\texttt{MACRODR\_\allowbreak AXIS\_\allowbreak SERIAL=1} is
load-bearing: it serialises the internal grid-combination loop, without which memory
grows with the number of concurrent combinations and the jobs exhaust memory.

\section*{S1.4\quad The optimiser schedule and its tolerances}

The script sets the damping schedule (initial damping $10$, growth
factor $3$, maximum damping $10^{10}$) and three limits: at most $100$ iterations, a
relative gradient tolerance of $10^{-6}$ and a value tolerance of $10^{-10}$.
The Newton-decrement criterion that stops most fits is described in Methods.

\end{document}
